\documentclass[aspectratio=169]{beamer}
\usetheme{Madrid}
\usecolortheme{default}
\usefonttheme{professionalfonts}

\usepackage[utf8]{inputenc}
\usepackage[T1]{fontenc}
\usepackage[spanish,es-nodecimaldot]{babel}
\usepackage{lmodern}
\usepackage{ragged2e}
\usepackage{booktabs}
\usepackage{enumitem}
\usepackage{hyperref}

\setbeamertemplate{navigation symbols}{}
\setbeamertemplate{footline}[frame number]
\setbeamertemplate{bibliography item}[text]

\title{Propuestas de plan de tesis\\Problema de investigación}
\subtitle{Curso: Proyecto de Investigación I - Maestría en Inteligencia Artificial}
\author{Luis Koc}
\date{\today}

\newcommand{\smallrefs}[1]{{\footnotesize #1}}

\begin{document}

\begin{frame}
  \titlepage
\end{frame}

\begin{frame}{Propósito de la presentación}
\justifying
Se presentan tres propuestas de plan de tesis elaboradas con un criterio común: priorizar la \textbf{comprensión del problema de investigación} antes de discutir metodologías, herramientas o recursos.

\vspace{0.5em}
En coherencia con la orientación del curso, el énfasis se coloca en:
\begin{itemize}[leftmargin=*]
    \item la realidad problemática,
    \item la formulación preliminar del problema,
    \item las preguntas de investigación,
    \item los objetivos centrados en comprender el problema, y
    \item los antecedentes que describen, cuantifican o discuten la problemática real.
    \item En esta etapa \textbf{no} se propone todavía una metodología de solución ni una arquitectura de inteligencia artificial.
\end{itemize}
\end{frame}

\section{Riesgos de ejecución}

\begin{frame}{Perfil del tesista y viabilidad de ejecución}
\justifying
El perfil del tesista muestra una combinación de fortalezas para una tesis aplicada: experiencia profesional extensa en energía, conociminetos del idioma inglés, capacidad de programación y familiaridad con herramientas digitales.

\vspace{0.5em}
No obstante, precisamente por tratarse de un perfil ejecutivo-técnico, la realización de la tesis enfrenta riesgos específicos que deben ser reconocidos desde el inicio como parte de la planificación académica.
\end{frame}

\begin{frame}{Riesgos principales de ejecución de la tesis}
\justifying
\begin{itemize}[leftmargin=*]
    \item \textbf{Restricción de tiempo y continuidad:} cargo de responsabilidad a tiempo completo.
    \item \textbf{Fragmentación atencional:} múltiples frentes operativos pueden reducir la concentración sostenida para desarrollar la investigación.
    \item \textbf{Ampliación excesiva del alcance:} la experiencia amplia puede llevar a formular un tema demasiado ambicioso para una maestría.
    \item \textbf{Sesgo hacia la solución prematura:} por el dominio técnico, existe tendencia a pensar antes en la herramienta que en el problema.
    \item \textbf{Riesgo de confidencialidad y acceso a datos:} algunos temas podrían depender de información empresarial no publicable.
    \item \textbf{Ajuste al registro académico:} la escritura profesional no siempre coincide con la escritura exigida en una tesis.
\end{itemize}
\end{frame}

\begin{frame}{Factores mitigantes y controles recomendados}
\justifying
El perfil también contiene factores mitigantes relevantes: disciplina profesional, madurez personal, base analítica sólida y capacidad para revisar literatura internacional.

\vspace{0.5em}
Se recomienda:
\begin{itemize}[leftmargin=*]
    \item reservar bloques horarios fijos para la tesis,
    \item mantener una sola pregunta central de investigación,
    \item separar la etapa de comprensión del problema de la etapa metodológica,
    \item elegir un tema sostenible con datos públicos o anonimizable, y
    \item trabajar con hitos mensuales breves y acumulativos.
\end{itemize}

\vspace{0.5em}
\textbf{Conclusión:} la tesis es claramente viable para este perfil, siempre que se gestione con rigor el tiempo disponible y, sobre todo, la delimitación del alcance.
\end{frame}

\section{Tema 1}

\begin{frame}{Tema 1: Título tentativo}
\textbf{{\Large Caracterización del problema térmico en cables eléctricos enterrados bajo condiciones reales de instalación y operación}}

\vspace{0.8em}
\justifying
Idea central: antes de modelar o predecir, es necesario comprender por qué el comportamiento térmico de un cable enterrado resulta incierto, variable y sensible a condiciones reales de suelo, ambiente e instalación.
\end{frame}

\begin{frame}{Tema 1: Realidad problemática}
\justifying
Los cables eléctricos enterrados son una infraestructura crítica en sistemas de transmisión, distribución y en plantas renovables eólicas y solares. Sin embargo, su capacidad de conducción depende de un fenómeno térmico complejo: la disipación de calor desde el conductor hacia el medio circundante.

\vspace{0.5em}
Ese proceso no depende solo del cable, sino también de variables del entorno:
\begin{itemize}[leftmargin=*]
    \item temperatura del suelo y ambiente,
    \item humedad,
    \item resistividad térmica del terreno,
    \item profundidad de enterramiento,
    \item material de relleno, e
    \item interacción térmica entre cables próximos u otras fuentes de calor.
\end{itemize}
\end{frame}

\begin{frame}{Tema 1: Núcleo del problema}
\justifying
En condiciones reales, estas variables rara vez son uniformes o estables. Cambian con el clima, la estación, la composición del suelo y la historia de carga. Cuando esta variabilidad no se comprende bien:
\begin{itemize}[leftmargin=*]
    \item la ampacidad puede sobrestimarse, elevando el riesgo de sobrecalentamiento y deterioro del aislamiento; o
    \item puede subestimarse, generando sobredimensionamiento y uso ineficiente de la infraestructura.
\end{itemize}

\vspace{0.5em}
\textbf{Problema preliminar:} no se comprende con suficiente precisión cómo las condiciones reales de instalación y del entorno térmico modifican la temperatura de operación y la capacidad de conducción de corriente de los cables eléctricos enterrados.

El diseño y dimensionamiento de los cables en la industria se basa en factores de corrección según normas conservadoras que no toman en cuenta casos de instlaciones complejas y la variación de las condiciones del entorno.
\end{frame}

\begin{frame}{Tema 1: Pregunta y objetivos}
\textbf{Pregunta general}

¿De qué manera las condiciones reales del suelo y del entorno de instalación inciden en el comportamiento térmico y en la capacidad de conducción de corriente de los cables eléctricos enterrados?

\vspace{0.7em}
\textbf{Objetivo general}

Analizar la naturaleza del problema térmico en cables eléctricos enterrados, identificando las condiciones físicas y de instalación que determinan su temperatura de operación y su capacidad de conducción de corriente.

\vspace{0.7em}
\textbf{Objetivos específicos}
\begin{itemize}[leftmargin=*]
    \item Describir las variables térmicas y geotécnicas relevantes.
    \item Examinar la influencia del entorno térmico real sobre la ampacidad.
    \item Identificar supuestos simplificadores de la literatura clásica.
    \item Delimitar vacíos de conocimiento en escenarios no ideales.
\end{itemize}
\end{frame}

\begin{frame}{Tema 1: Antecedentes centrados en el problema}
\justifying
Los antecedentes muestran que el problema térmico depende del entorno real y no puede reducirse a parámetros nominales del cable:
\begin{itemize}[leftmargin=*]
    \item \textbf{Aras, Oysu y Yilmaz (2005):} evidencian limitaciones de los métodos convencionales frente a configuraciones complejas.
    \item \textbf{Ocłoń et al. (2015):} muestran que el espaciamiento, la geometría del lecho térmico y el material de relleno alteran la temperatura del núcleo.
    \item \textbf{Kim et al. (2025):} señalan que temperatura ambiente y del suelo modifican la temperatura máxima y la ampacidad crítica.
    \item \textbf{Khumalo et al. (2025):} revelan que secado del suelo y cambios de resistividad térmica reducen significativamente la ampacidad.
\end{itemize}
\end{frame}

\section{Tema 2}

\begin{frame}{Tema 2: Título tentativo}
\textbf{{\Large Caracterización del problema de incertidumbre en el pronóstico con 24 horas de anticipación de la producción de energía renovable}}

\vspace{0.8em}
\justifying
La producción electrica de las plantas de energía renovable es variable (solar, eólica) y es difícil de anticipar de manera confiable en un horizonte operativo de 24 horas.
\end{frame}

\begin{frame}{Tema 2: Realidad problemática}
\justifying
La expansión de la generación renovable variable ha intensificado la necesidad de anticipar la producción esperada para el día siguiente. El horizonte de 24 horas es particularmente importante porque coincide con decisiones de:
\begin{itemize}[leftmargin=*]
    \item programación de unidades,
    \item asignación de reservas,
    \item participación en mercados eléctricos, y
    \item coordinación de recursos flexibles.
\end{itemize}

\vspace{0.5em}
Sin embargo, la generación eólica y fotovoltaica depende de fenómenos meteorológicos variables e inherentemente inciertos.
\end{frame}

\begin{frame}{Tema 2: Núcleo del problema}
\justifying
El error de pronóstico no es simplemente ruido estadístico. Expresa una dificultad física y operativa asociada a:
\begin{itemize}[leftmargin=*]
    \item nubosidad,
    \item viento,
    \item irradiancia,
    \item temperatura,
    \item heterogeneidad espacial del recurso, y
    \item calidad de los pronósticos meteorológicos de entrada.
\end{itemize}

\vspace{0.5em}
\textbf{Problema preliminar:} la producción de energía renovable presenta un nivel de variabilidad e incertidumbre que dificulta anticipar con suficiente confiabilidad, a 24 horas, la energía que será realmente entregada por plantas eólicas o fotovoltaicas.
\end{frame}

\begin{frame}{Tema 2: Pregunta y objetivos}
\textbf{Pregunta general}

¿Cuáles son los factores que explican la incertidumbre del pronóstico con 24 horas de anticipación de la producción de energía renovable y cómo se expresa dicha incertidumbre en la operación eléctrica?

\vspace{0.7em}
\textbf{Objetivo general}

Analizar la naturaleza del problema de incertidumbre asociado al pronóstico day-ahead de producción renovable, identificando los factores físicos, meteorológicos y operativos que dificultan su representación confiable.

\vspace{0.7em}
\textbf{Objetivos específicos}
\begin{itemize}[leftmargin=*]
    \item Describir las fuentes de variabilidad e incertidumbre del recurso renovable.
    \item Examinar por qué el horizonte de 24 horas es operativamente crítico.
    \item Identificar consecuencias técnicas y económicas del error de pronóstico.
    \item Delimitar vacíos de conocimiento sobre la incertidumbre day-ahead.
\end{itemize}
\end{frame}

\begin{frame}{Tema 2: Antecedentes centrados en el problema}
\justifying
Los antecedentes resaltan la naturaleza estructural de la incertidumbre del pronóstico renovable:
\begin{itemize}[leftmargin=*]
    \item \textbf{Antonanzas et al. (2016):} muestran que la variabilidad solar impacta la gestión de red, el despacho y la estabilidad del sistema.
    \item \textbf{Hanifi et al. (2020):} identifican la incertidumbre y fluctuación del viento como obstáculo central para la integración eólica.
    \item \textbf{Leva et al. (2017):} subrayan la relevancia operativa específica del horizonte de 24 horas para redes inteligentes y mercados day-ahead.
\end{itemize}
\end{frame}

\section{Tema 3}

\begin{frame}{Tema 3: Título tentativo}
\textbf{{\Large Identificación del problema de desconocimiento del estado térmico real y de los cuellos de botella térmicos en cables subterráneos de potencia}}

\vspace{0.8em}
\justifying
Esta propuesta se concentra en una pregunta previa a cualquier modelo predictivo: por qué el estado térmico real de un cable subterráneo es difícil de observar, localizar e interpretar a lo largo de su ruta.
\end{frame}

\begin{frame}{Tema 3: Realidad problemática}
\justifying
Los cables subterráneos suelen operar con criterios de capacidad térmica estática definidos en la etapa de diseño. Tales criterios son seguros, pero normalmente descansan sobre supuestos conservadores de ambiente, suelo y condiciones de instalación.

\vspace{0.5em}
En operación real, el entorno del cable es heterogéneo y puede generar tramos con menor eficiencia de disipación de calor. Allí aparecen \textbf{cuellos de botella térmicos} que terminan limitando la capacidad del sistema completo.
\end{frame}

\begin{frame}{Tema 3: Núcleo del problema}
\justifying
La dificultad no reside solo en que existan puntos térmicamente críticos, sino en que el estado térmico real del conductor no siempre es directamente observable. El acceso físico es limitado y esto dificulta conocer con precisión:
\begin{itemize}[leftmargin=*]
    \item la temperatura real del cable,
    \item la distribución de calor a lo largo del tendido, y
    \item el margen térmico disponible en cada tramo.
\end{itemize}

\vspace{0.5em}
\textbf{Problema preliminar:} existe un conocimiento insuficiente sobre la ubicación y el comportamiento de los cuellos de botella térmicos en cables subterráneos, así como sobre la diferencia entre la capacidad térmica nominal y la capacidad térmica realmente disponible en condiciones operativas cambiantes.
\end{frame}

\begin{frame}{Tema 3: Pregunta y objetivos}
\textbf{Pregunta general}

¿Cómo se caracteriza el estado térmico real y de los cuellos de botella térmicos en cables subterráneos de potencia bajo condiciones operativas y ambientales variables?

\vspace{0.7em}
\textbf{Objetivo general}

Analizar el problema del desconocimiento del estado térmico real en cables subterráneos de potencia, con especial atención a la formación de cuellos de botella térmicos y a la diferencia entre capacidad estática y capacidad efectiva de carga.

\vspace{0.7em}
\textbf{Objetivos específicos}
\begin{itemize}[leftmargin=*]
    \item Describir condiciones que favorecen zonas térmicamente críticas.
    \item Examinar la relación entre heterogeneidad del entorno y capacidad de carga.
    \item Identificar vacíos de información relevantes para la gestión térmica.
    \item Sintetizar evidencia sobre puntos calientes y cuellos de botella.
\end{itemize}
\end{frame}

\begin{frame}{Tema 3: Antecedentes centrados en el problema}
\justifying
Los antecedentes muestran que la capacidad térmica efectiva depende de información localizada que a menudo no está disponible:
\begin{itemize}[leftmargin=*]
    \item \textbf{Li, Tan y Su (2006):} plantean que la evaluación en tiempo real exige identificar puntos calientes y parámetros térmicos del suelo a lo largo del cable.
    \item \textbf{Lai y Teh (2022):} muestran que los límites térmicos de cables pueden convertirse en cuellos de botella no bien caracterizados por enfoques estáticos.
    \item \textbf{Kim et al. (2025):} refuerzan que la capacidad disponible depende del entorno real.
    \item \textbf{Khumalo et al. (2025):} evidencian que humedad, resistividad térmica y profundidad alteran fuertemente la ampacidad.
\end{itemize}
\end{frame}

\section{Cierre}

\begin{frame}{Cierre comparativo}
\justifying
Las tres propuestas comparten un mismo criterio epistemológico: \textbf{partir de un problema real antes de discutir la solución}.

\vspace{0.5em}
\begin{itemize}[leftmargin=*]
    \item \textbf{Tema 1:} problema físico del comportamiento térmico de cables enterrados.
    \item \textbf{Tema 2:} problema operativo de la incertidumbre day-ahead de la generación renovable.
    \item \textbf{Tema 3:} problema de visibilidad y caracterización térmica de activos subterráneos en operación.
\end{itemize}

\vspace{0.5em}
Cualquiera de las tres propuestas puede evolucionar después hacia una tesis en inteligencia artificial, pero en la etapa actual su valor principal está en delimitar con claridad la realidad problemática, la pregunta central y los antecedentes que justifican investigar.
\end{frame}

\begin{frame}[allowframebreaks]{Referencias}
\footnotesize
\justifying
Aras, F., Oysu, C., \& Yilmaz, G. (2005). An assessment of the methods for calculating ampacity of underground power cables. \textit{Electric Power Components and Systems, 33}(12), 1385--1402. \url{https://doi.org/10.1080/15325000590964425}.
\bigskip

Antonanzas, J., Osorio, N., Escobar, R., Urraca, R., Martinez-de-Pison, F. J., \& Antonanzas-Torres, F. (2016). Review of photovoltaic power forecasting. \textit{Solar Energy, 136}, 78--111. \url{https://doi.org/10.1016/j.solener.2016.06.069}.
\bigskip

Hanifi, S., Liu, X., Lin, Z., \& Lotfian, S. (2020). A critical review of wind power forecasting methods--Past, present and future. \textit{Energies, 13}(15), 3764. \url{https://doi.org/10.3390/en13153764}.
\bigskip

Khumalo, N. Q., Naidoo, R. M., Mbungu, N. T., \& Bansal, R. C. (2025). A critical assessment of cable rating methods under soil drying out conditions. \textit{International Transactions on Electrical Energy Systems, 2025}, Article 5946564. \url{https://doi.org/10.1155/etep/5946564}.
\bigskip

Kim, Y.-S., Nguyen Cong, H., Dinh, B. H., \& Kim, H.-K. (2025). Effect of ambient air and ground temperatures on heat transfer in underground power cable system buried in newly developed cable bedding material. \textit{Geothermics, 125}, 103151. \url{https://doi.org/10.1016/j.geothermics.2024.103151}.
\bigskip

Lai, C.-M., \& Teh, J. (2022). Comprehensive review of the dynamic thermal rating system for sustainable electrical power systems. \textit{Energy Reports, 8}, 3263--3288. \url{https://doi.org/10.1016/j.egyr.2022.02.085}.
\bigskip

Leva, S., Dolara, A., Grimaccia, F., Mussetta, M., \& Ogliari, E. (2017). Analysis and validation of 24 hours ahead neural network forecasting of photovoltaic output power. \textit{Mathematics and Computers in Simulation, 131}, 88--100. \url{https://doi.org/10.1016/j.matcom.2015.05.010}.
\bigskip

Li, H.-J., Tan, K. C., \& Su, Q. (2006). Assessment of underground cable ratings based on distributed temperature sensing. \textit{IEEE Transactions on Power Delivery, 21}(4), 1763--1769. \url{https://doi.org/10.1109/TPWRD.2006.874668}.
\bigskip

Ocłoń, P., Cisek, P., Taler, D., Pilarczyk, M., \& Szwarc, T. (2015). Optimizing of the underground power cable bedding using momentum-type particle swarm optimization method. \textit{Energy, 92}(P2), 230--239. \url{https://doi.org/10.1016/j.energy.2015.04.100}.
\end{frame}

\begin{frame}
\centering
\Large Gracias
\end{frame}

\end{document}
